\section{Requerimientos}

\begin{enumerate}
    \item \textbf{RF01-Módulo 1:Sistema de Sensores}
        \begin{enumerate}
            \item \textbf{Verificaciones de validez de las entradas}
                \begin{enumerate}
                    \item El sensor de pH deberá ser capaz de trabajar en líquidos en un rango de temperatura de 18°C a 30°C, sin contaminar ni interferir con la fermentación del kéfir, atraves de una precisión de ±0.1 pH ,un tiempo de respuesta de 1 minuto y un rango de medición de 4.0 a 7.0 de pH.
                    \item El sensor de temperatura deberá ser capaz de trabajar en líquidos , sin contaminar ni interferir con la fermentación del kéfir, atraves de una precisión de ±0.5 °C y un rango de medición de 18°C a 30°C.
                    \item El sistema deberá contar con un sensor de $CO_{2}$  capaz de medir niveles en el aire dentro del frasco de fermentación, operando en el mismo rango de temperatura (18°C a 30°C).
                    \item El microcontrolador debera tener al menos 3 entradas analógicas para conectar los sensores de pH, temperatura y $CO_{2}$.
                    \item El microcontrolador deberá tener una memoria suficiente para almacenar datos temporalmente en caso de fallos de conexión con la base de datos.
                    \item El microcontroladror debera permitir enviar los datos de manera sincróno a la base de datos.
                \end{enumerate}   
            
            \item \textbf{Secuencia exacta de operaciones}
                \begin{enumerate}
                    \item El sistema deberá contar con dos modos de operación: modo prueba y modo monitoreo, activados mediante dos botones físicos en el microcontrolador que, al presionarse, inicien el modo correspondiente y permitan salir del mismo al presionar nuevamente el botón.
                    \item En modo prueba, el sistema deberá verificar el estado de los sensores al enviar los datos de los sensores cada 1 minuto exclusivamente a la página web para visualización temporal, sin almacenarlos en la base de datos.
                    \item En modo monitoreo, el sistema deberá verificar que el pH inicial de la leche esté entre 6.6 a 6.8 y su temperaturaesté entre 18°C a 30°C, en caso contrario, no iniciar el monitoreo y enviar una alerta a la base de datos.
                    \item En modo monitoreo, el sistema deberá verificar la conexión con la base de datos antes de iniciar la recolección de datos, y en caso de fallo, no iniciar el monitoreo y enviar una alerta.
                    \item En modo monitoreo, el sistema deberá iniciar la recolección y envío de datos a la base de datos.
                    \item En modo monitoreo, el sistema deberá enviar los datos recopilados por los sensores de pH y $CO_{2}$  cada 10 minutos si el pH está en el rango de 6.8 a 6.0.
                    \item En modo monitoreo, el sistema deberá enviar los datos recopilados por los sensores cada de pH y $CO_{2}$  cada 5 minutos si el pH está en el rango de 6.0 a 5.0.
                    \item En modo monitoreo, el sistema deberá recopilar y enviar los datos de los tres sensores (pH, $CO_{2}$  y temperatura) cada 1 minuto si el pH está en el rango de 5.0 a 4.2 en adelante.
                    \item En modo monitoreo, el sistema deberá recopilar y enviar los datos del sensor de temperatura cada 2.5 minutos de manera independiente, y ajustar la frecuencia a cada 1 minuto de todos los sensores si la temperatura supera los 28°C.
                    \item El microcontrolador deberá recopilar los datos de los sensores de manera secuencial (primero pH, luego $CO_{2}$ , luego temperatura) en cada ciclo de medición para optimizar el procesamiento.
                \end{enumerate}   
        
            \item \textbf{Respuestas a situaciones anormales}
                \begin{enumerate}
                    \item El sistema deberá generar una alerta y registrarla en la base de datos si la comunicación con cualquier sensor (pH, temperatura, CO) se interrumpe ya sea en modo prueba o monitoreo.
                    \item El sistema deberá reenviar los datos recolectados por los 3 sensores en caso de un fallo en la conexión con la base de datos.
                    \item El sistema deberá registrar en un log todos los intentos de conexión fallidos al servidor de base de datos.
                \end{enumerate}   
       
            \item \textbf{Efecto de los parámetros}
                \begin{enumerate} 
                    \item El microcontrolador deberá ajustar la frecuencia de recolección de datos en función del rango de pH medido y la temperaura, como se especifica en la secuencia de operaciones.
                    \item La alimentación del microcontrolador deberá ser estable y adecuada para garantizar el funcionamiento continuo durante el proceso de fermentación, considerando al menos un intervalo de 12 a 48 horas de operación ininterrumpida.
                \end{enumerate}   
        
            \item \textbf{Relación entre las salidas y las entradas}
                \begin{enumerate}
                    \item El microcontrolador debera recibir lecturas crudas(haciendo la conversión de datos analogicos a digitales) de los sensores y genera paquetes de datos estructurados junto con un identificador del microcontrolador para su envío a la base de datos.
                    \item El protocolo de conexión del microcontrolador debera ser compatible con el servidor de base de datos para asegurar la correcta transmisión de datos.
                    \item El protocolo de conexión deberá garantizar que los datos enviados a la base de datos incluyan una marca de tiempo para cada lectura, permitiendo un seguimiento preciso de la evolución de los parámetros monitoreados.
                    \item El protocolo de conexión deberá asegurar que los datos enviados a la base de datos estén en el formato requerido por el módulo de análisis de datos para su procesamiento posterior.
                \end{enumerate}
        \end{enumerate}

    \item \textbf{Módulo 2:Analisis de Datos} 
        \begin{enumerate}
            \item \textbf{Verificaciones de validez de las entradas}
                \begin{enumerate}
                    \item El sistema deberá verificar que los datos recibidos de la base de datos estén completos y no contengan valores nulos.
                    \item El sistema deberá verificar que los datos de pH estén dentro del rango esperado (4.0 a 7.0).
                    \item El sistema deberá verificar que los datos de temperatura estén dentro del rango esperado (0°C a 40°C).
                    \item El sistema deberá verificar que los datos de CO estén dentro del rango esperado (0 a 1000 ppm).
                    El sistema deberá verificar que los datos incluyan un timestamp válido y estén ordenados cronológicamente para garantizar la correcta secuencia temporal del análisis.
                \end{enumerate}   
            
            \item \textbf{Secuencia exacta de operaciones}
                \begin{enumerate}
                    \item El sistema deberá iniciar el análisis de datos automáticamente cada hora.
                    \item El sistema deberá procesar los datos en el orden en que fueron recibidos.
                    \item El sistema deberá generar un informe de análisis después de procesar un lote completo de datos.
                \end{enumerate}   
        
            \item \textbf{Respuestas a situaciones anormales}
                \begin{enumerate}
                    \item El sistema deberá generar una alerta si se detectan valores fuera del rango esperado en cualquiera de los parámetros (pH, temperatura, CO).
                    \item El sistema deberá registrar en un log cualquier error encontrado durante el análisis de datos.
                    \item El sistema deberá notificar al administrador si se detectan patrones anormales en los datos, como una caída rápida del pH o un aumento inusual en los niveles de CO.
                \end{enumerate}   
       
            \item \textbf{Efecto de los parámetros}
                \begin{enumerate}
                    \item El sistema deberá ajustar la frecuencia del análisis de datos si se detectan cambios significativos en los parámetros monitoreados.
                    \item El sistema deberá priorizar el análisis de datos si se detectan valores críticos en cualquiera de los parámetros.
                \end{enumerate}   
        
            \item \textbf{Relación entre las salidas y las entradas}
                \begin{enumerate}
                    \item El sistema deberá generar informes detallados basados en los datos analizados, incluyendo gráficos y tendencias.
                    \item El sistema deberá actualizar la base de datos con los resultados del análisis para su posterior consulta.
               \end{enumerate}
        \end{enumerate}
    \item \textbf{Módulo 3:Base de Datos}
        \begin{enumerate}
            \item \textbf{Verificaciones de validez de las entradas}
                \begin{enumerate}
                    \item El sistema deberá verificar que los datos recibidos del sistema de sensores estén completos y no contengan valores nulos.
                    \item El sistema deberá verificar que los datos de pH estén dentro del rango esperado (4.0 a 7.0).
                    \item El sistema deberá verificar que los datos de temperatura estén dentro del rango esperado (0°C a 40°C).
                    \item El sistema deberá verificar que los datos de CO estén dentro del rango esperado (0 a 1000 ppm).
                \end{enumerate}   
            
            \item \textbf{Secuencia exacta de operaciones}
                \begin{enumerate}
                    \item El sistema deberá almacenar los datos recibidos en la base de datos en el orden en que fueron recibidos.
                    \item El sistema deberá realizar copias de seguridad automáticas de la base de datos cada 24 horas.
                    \item El sistema deberá permitir consultas y recuperaciones de datos por parte del módulo de análisis de datos.
                \end{enumerate}   
        
            \item \textbf{Respuestas a situaciones anormales}
                \begin{enumerate}
                    \item El sistema deberá generar una alerta si se detectan intentos no autorizados de acceso a la base de datos.
                    \item El sistema deberá registrar en un log cualquier error encontrado durante las operaciones de almacenamiento o recuperación de datos.
                    \item El sistema deberá notificar al administrador si se detectan problemas recurrentes con la integridad de los datos.
                \end{enumerate}   
       
            \item \textbf{Efecto de los parámetros}
                \begin{enumerate}
                    \item El sistema deberá ajustar la frecuencia de las copias de seguridad si se detectan cambios significativos en el volumen de datos almacenados.
                    \item El sistema deberá optimizar las consultas si se detectan patrones comunes en las solicitudes de datos.
                \end{enumerate}   
        
            \item \textbf{Relación entre las salidas y las entradas}
                \begin{enumerate}
                    \item El sistema deberá proporcionar acceso rápido y eficiente a los datos almacenados para su análisis y visualización.
                    \item El sistema deberá mantener la integridad y seguridad de los datos almacenados en la base de datos.
                \end{enumerate}
        \end{enumerate}

    \item \textbf{Módulo 4:Página web} 
        \begin{enumerate}
            \item \textbf{Verificaciones de validez de las entradas}
                \begin{enumerate}
                    \item El sistema deberá verificar que los datos recibidos del módulo de análisis de datos estén completos y no contengan valores nulos.
                    \item El sistema deberá verificar que los datos de pH, temperatura y CO estén dentro del rango esperado antes de mostrarlos en la interfaz web.
                \end{enumerate}   
            
            \item \textbf{Secuencia exacta de operaciones}
                \begin{enumerate}
                    \item El sistema deberá actualizar la interfaz web con los datos más recientes cada 10 minutos.
                    \item El sistema deberá permitir a los usuarios filtrar y ordenar los datos mostrados en la interfaz web.
                    \item El sistema deberá proporcionar gráficos y visualizaciones basadas en los datos analizados.
                \end{enumerate}   
        
            \item \textbf{Respuestas a situaciones anormales}
                \begin{enumerate}
                    \item El sistema deberá mostrar un mensaje de error si no se pueden cargar los datos debido a problemas de conexión con la base de datos.
                    \item El sistema deberá registrar en un log cualquier error encontrado durante la actualización o visualización de datos en la interfaz web.
                    \item El sistema deberá notificar al administrador si se detectan problemas recurrentes con la carga de datos en la interfaz web.
                \end{enumerate}   
       
            \item \textbf{Efecto de los parámetros}
                \begin{enumerate}
                    \item El sistema deberá ajustar la frecuencia de actualización de la interfaz web si se detectan cambios significativos en los parámetros monitoreados.
                    \item El sistema deberá optimizar la visualización de datos si se detectan patrones comunes en las solicitudes de los usuarios.
                \end{enumerate}   
        
            \item \textbf{Relación entre las salidas y las entradas}
                \begin{enumerate}
                    \item El sistema deberá proporcionar una interfaz intuitiva y fácil de usar para que los usuarios puedan acceder y comprender los datos monitoreados.
                    \item El sistema deberá garantizar que los datos mostrados en la interfaz web sean precisos y reflejen el estado actual del monitoreo.
                \end{enumerate}
        \end{enumerate}

\end{enumerate}
